	\documentclass[12pt]{article}

		\usepackage{graphicx}
		\usepackage{amsmath,amssymb,rotating,multirow,setspace,fullpage,verbatim,subfigure,quotes,url,bm}

	\usepackage[left=1in,right=1in,top=1in,bottom=1.1in]{geometry}

		\begin{document}

			
				\begin{center}
				{\large \textbf{QTM 220: Regression Analysis } }\\
				\vspace{0.1in}
				Emory University\\
				\vspace{0.2in}
				\begin{tabular}{llll}
				&	Lecture room:     & PAIS 230  &\\
					& Lecture time:   & Mondays and Wednesdays, 2:30-3:45pm & \\
					& Lab room:       & PAIS 230 & \\
				& 	Lab time: & Fridays, 2:30-3:20pm & \\
					\\
					& Instructor:     & David A. Hirshberg & \\
					&Email: & davidahirshberg@emory.edu &\\
					&Office: & PAIS 583  &\\
					& Office hours: TBD &  \\
					& TA: & TBD &\\
					& Email: & & \\
					
				\end{tabular}
			\end{center}
			
			
\subsection*{Course Description}
This course introduces students to widely-used procedures for regression analysis and provides intuitive, applied, and formal foundations for more advanced methods treated in later courses. The course covers descriptive, population and causal inference with regression from a modern perspective. While the course will emphasize the mathematical foundations of these concepts, each topic will also cover the implementation of the relevant methods in an application through the use of the statistical programming language R.

This course represents the culmination of the required courses for QTM majors. It uses the coding skills learned in QTM 150 and the math skill learned in Math 210 and QTM 210 to conduct much of the analysis discussed in QTM 110. QTM 220 also provides an introduction and a foundation for future study of topics such as machine learning and causal inference. 


	
\bigskip


\paragraph{Course Goals} By the end of this course, students will be able to 


\begin{itemize}
\item Formalize questions about the world, differentiating between questions about a sample, a population from which we've sampled, and counterfactual quantities including causal effects.
\item Formalize random sampling and randomization assumptions that allow those questions to be answered using the sample we've observed.
\item Use regression to answer those questions.
\item Quantify uncertainty and error.
\item Understand when certain techniques are appropriate to use (or not).
\item Demonstrate proficiency in statistics-relevant programming skills in the R language, including data simulation, visualization, and modeling.
\end{itemize}

\paragraph{Prerequisites}


\subsubsection*{Course Requirements} 
			\paragraph{Attendance} Lecture attendance is expected. If a student cannot attend lecture, they must email the instructor to explain their absence. Please do not come to class (lecture or lab) if sick. 
			
			\paragraph{Problem Sets} Problem sets will be assigned approximately weekly throughout the semester. The homework assignments will consist of analytical problems, computer work, and/or data analysis.
			Assignments will be posted on and submitted via Canvas. No late work will be accepted. We encourage students to confer with each other on the homework assignments, but you must write your own solutions.
			The use of Large Language Models, e.g. GPT4, to assist you in writing them is permitted. However, you are expected to be able to explain anything you turn in. 
			
			\paragraph{Midterm Exams} The midterms will be in-class exams {\it  There is no collaboration allowed on the midterms.}

			\paragraph{Final Exam}  The final will be held during the assigned exam time and location. {\it  There is no collaboration allowed on the final.}  
			
			\paragraph{Lab Session} Lab sessions will be run by the TAs. The purpose of the lab sessions is to learn the R programming and data analysis skills useful for regression and other empirical analysis. Students must bring a laptop, with R installed and updated on it, to these sessions. Lab sessions will prepare students to complete the problem sets and review the solutions of previous homeworks. Lab participation and course performance are highly correlated. 
			
%			\paragraph*{Policies} any other course policies??? laptops in lecture? using Canvas discussion threads to ask questions about course material instead of emailing? extra credit (and lack thereof)? grade appeals? Pablo's grade appeal policy: If you believe that your grade on any assignment or exam question is incorrect or unfair, you should submit your concerns, in writing, to the instructor.  There are two types of appeals: 
%1.	Addition errors or other administrative snafus: Please let me know if you catch one of these ASAP.  I’ll confirm the error and change your grade in the gradebook. 
%2.	Disagreement about grading criteria:  Please wait two days (or the nearest non-weekend day) after receiving your grade before submitting a written appeal.  Do not leave comments on Canvas assignments for TAs as they are not responsible for grade appeals or score corrections.  The written appeal should fully summarize what you believe the problems are and why.   The instructor will consider your appeal. If you are not satisfied with the response, you may resubmit the assignment for regrading in its entirety by the instructor  This grade will be final.  Note that grades may go up or down during an appeal.  Failure to comply with this appeal process means that you forfeit your appeal. 

		
	\subsubsection*{Grading }		
	
			\begin{itemize}
			\item Problem Sets (35\%)
			\item Midterm Exams (40\%)
			\item Final Exam (25\%)
			\end{itemize}
			
			\subsubsection*{Incomplete Grades}
			No incomplete grades will be given unless there is an agreement between the instructor and the student {\bf prior} to the end of the course.  The instructor retains the right to determine legitimate reasons for an incomplete grade.
			
\subsubsection*{Honor Code}

All students enrolled at Emory are expected to abide by the Emory College Honor Code. Any type of academic misconduct is not allowed which includes 1) receiving or giving information about the content or conduct of an examination knowing that the release of such information is not allowed and 2) plagiarizing, whether intentionally or unintentionally, in any assignment. For the activities that are considered to be academically dishonest, refer to the Honor Code: http://catalog.college.emory.edu/academic/policies-regulations/honor-code.html.\\

\subsubsection*{Accessibility and Accommodations}

If you are seeking classroom accommodations or academic adjustments under the Americans with Disabilities Act, you are required to register with Office of Accessibility Services (http://accessibility.emory.edu). To receive academic accommodations for this class, please obtain the relevant letter and meet with the instructor at the beginning of the semester. Students are expected to give two weeks notice of the need for accommodations.
%Emory University is committed under the Americans with Disabilities Act and its Amendments and Section 504 of the Rehabilitation Act to providing appropriate accommodations to individuals with documented disabilities. If you have a disability-related need for reasonable academic adjustments in this course, provide the instructor(s) with an accommodation notification letter from Access, Disabilities Services and Resources office (\href{http://www.ods.emory.edu/}{http://www.ods.emory.edu/}). Students are expected to give two weeks notice of the need for accommodations. If you need immediate accommodations or physical access, please arrange to meet with instructor as soon as your accommodations have been finalized.\\

\bigskip
			\newpage 
\section*{Tentative Schedule}
\subsection*{Part 1: One and Two-Sample Problems}

\subsubsection*{Lecture 1 (Jan 17)}
Summarizing Samples.

\subsubsection*{Lecture 2 (Jan 22)}
Random Sampling and Sampling Distributions. Application: Polling to predict election turnout.

\subsubsection*{Lecture 3 (Jan 24)}
Interval Estimation and Confidence Intervals. Application: Polling to predict election turnout.

\subsubsection*{Lecture 4 (Jan 29)}
Estimating Sampling Distributions using the Bootstrap and Normal Approximation. Application: Polling to predict election turnout.

\subsubsection*{Lecture 5 (Jan 31)}
Working with Non-Binary Outcomes. Application: Understanding the relationship between education and income.

\subsubsection*{Lecture 6 (Feb 5)}
Two-Sample Problems, Part I. Application: Polling to predict election turnout.

\subsubsection*{Lecture 7 (Feb 7)}
Two-Sample Problems, Part II. Application: Understanding the relationship between education and income.

\subsubsection*{Lecture 8 (Feb 12)}
Causal Inference. Application: Estimating the effect of a work-training program on subsequent income.

\subsubsection*{Lecture 9 (Feb 14)}
Other things you may have heard of. Regression lines, statistical significance, one and two-sided tests, and homoskedasticity.

\subsubsection*{Review (Feb 19)}
\subsubsection*{Exam (Feb 21)}

\subsection*{Part 2: Regression Analysis with Multiple and Many-Valued Covariates}

\subsubsection*{Lecture 10 (Feb 26)}
Summarizing Samples with Many-valued Covariates.

\subsubsection*{Lecture 11 (Feb 28)}
Interval Estimation and Confidence Intervals with Many-valued Covariates.

\subsubsection*{Lecture 12 (Mar 4)}
Causal Inference with Many-valued Covariates.

\subsubsection*{Lecture 13 (Mar 6)}
Misspecification and Coverage.

\subsubsection*{Spring Break (Mar 11)}
\subsubsection*{Spring Break (Mar 13)}

\subsubsection*{Lecture 14 (Mar 18)}
Summarizing Samples using Bivariate Regression, Part I. Simpson's Paradox in Qualitative Terms.

\subsubsection*{Lecture 15 (Mar 20)}
Summarizing Samples using Bivariate Regression, Part II. Omitted Variable Bias and Simpson's Paradox in Quantitative Terms.

\subsubsection*{Lecture 16 (Mar 25)}
Multivariate Regression in General. Matrix/Vector Formulation.

\subsubsection*{Lecture 17 (Mar 27)}
Summarizing Populations using Multivariate Regression.

\subsubsection*{Lecture 18 (Apr 1)}
Causal Inference in Conditionally Randomized Experiments using Multivariate Regression. 

\subsubsection*{Review (Apr 3)}
\subsubsection*{Exam (Apr 8)}

\subsection*{Part 3: Trying To Prevent Bias From Telling the Story}

\subsubsection*{Lecture 19 (Apr 10)}
Avoiding Bias using Inverse Probability Weighting 
\subsubsection*{Lecture 20 (Apr 15)}
Model Selection via Cross-Validation         
\subsubsection*{Lecture 21 (Apr 17)}
Regression Trees  
\subsubsection*{Lecture 22 (Apr 22)}
Regression Forests 
\subsubsection*{TBD    (Apr 24)}
\subsubsection*{Review (Apr 29)}

\end{document} 
